% =====================================================================
%  FICHE 8 — THÈME 2 — NIVEAU FACILE — ÉNONCÉS
% =====================================================================
\documentclass[11pt,a4paper]{article}
\usepackage[utf8]{inputenc}
\usepackage[T1]{fontenc}
\usepackage{lmodern}
\usepackage[french]{babel}
\usepackage{amsmath,amssymb,mathtools}
\usepackage{icomma}
\usepackage{siunitx}
\sisetup{output-decimal-marker={,},per-mode=power,group-separator={\,},group-minimum-digits=5}
\usepackage[version=4]{mhchem}
\usepackage{xcolor}
\usepackage{tikz}
\usepackage[most]{tcolorbox}
\usepackage{enumitem}
\usepackage{array}
\usepackage{fancyhdr}
\usepackage[hidelinks]{hyperref}
\usetikzlibrary{arrows.meta,positioning,calc}
\usepackage[a4paper,margin=2.1cm,headheight=26pt,headsep=10pt]{geometry}
\usepackage{titlesec}

\definecolor{primary}{RGB}{150,98,162}
\definecolor{primarydark}{RGB}{92,52,110}

\tcbset{boxbase/.style={enhanced,breakable,boxrule=0.9pt,arc=2.5pt,left=3mm,right=3mm,top=2.3mm,bottom=2.3mm,fonttitle=\bfseries,coltitle=white,toptitle=0.6mm,bottomtitle=0.6mm}}
\newtcolorbox[auto counter]{exercice}[1][]{boxbase,colback=primary!4,colframe=primary,title={\textsc{Exercice}~\thetcbcounter\ \ifx\relax#1\relax\relax\else\textmd{--- \itshape #1}\fi}}

\newcommand{\dH}{\Delta H}
\newcommand{\dHr}{\Delta H^{\circ}_{\text{r}}}
\newcommand{\dHf}{\Delta H^{\circ}_{\text{f}}}

\pagestyle{fancy}
\fancyhf{}
\fancyhead[L]{\small\textcolor{primarydark}{\textbf{Fiche 8} — Thème 2 : Enthalpie de formation et loi de Hess}}
\fancyhead[R]{\small\textcolor{primarydark}{\textsc{Facile}}}
\fancyfoot[C]{\small\thepage}
\renewcommand{\headrulewidth}{0.8pt}
\renewcommand{\headrule}{\hbox to\headwidth{\color{primary}\leaders\hrule height \headrulewidth\hfill}}
\setlength{\parindent}{0pt}\setlength{\parskip}{3pt}

\begin{document}

\begin{center}
\begin{tikzpicture}
\node[rectangle,rounded corners=7pt,fill=primary,text=white,minimum width=16cm,
      minimum height=1.1cm,align=center,inner sep=6pt]
  {{\large\bfseries Thème 2 — Enthalpie de formation et loi de Hess}\\[1pt]{\small Niveau \textsc{Facile} — énoncés}};
\end{tikzpicture}
\end{center}

\begin{exercice}[reconnaître un $\dHf$ nul]
Parmi les espèces suivantes, indiquer celles dont l'enthalpie standard de formation $\dHf$ est
\textbf{nulle par convention}, et justifier en une phrase pour chacune.
\[ \ce{O2}(g),\quad \ce{CO2}(g),\quad \ce{Fe}(s),\quad \ce{H2O}(l),\quad \ce{N2}(g),\quad \ce{NH3}(g),
   \quad \ce{C}(\text{graphite}),\quad \ce{O3}(g). \]
\emph{(Indice pour l'ozone $\ce{O3}$ : la forme la plus stable de l'élément oxygène est $\ce{O2}$.)}
\end{exercice}

\begin{exercice}[écrire la loi de Hess]
Écrire l'expression de $\dHr$ \emph{en fonction des $\dHf$} pour chacune des réactions (sans valeur numérique) :
\begin{enumerate}[label=\alph*),leftmargin=2em,topsep=2pt,itemsep=2pt]
  \item $\ce{A} \rightarrow 2\,\ce{B}$ ;
  \item $\ce{N2}(g) + 3\,\ce{H2}(g) \rightarrow 2\,\ce{NH3}(g)$ ;
  \item $2\,\ce{C}(s) + 3\,\ce{H2}(g) + \tfrac12\,\ce{O2}(g) \rightarrow \ce{C2H5OH}(l)$.
\end{enumerate}
Simplifier chaque expression en tenant compte des corps purs simples.
\end{exercice}

\begin{exercice}[application numérique de la loi de Hess]
Pour la combustion du méthane \;$\ce{CH4}(g) + 2\,\ce{O2}(g) \rightarrow \ce{CO2}(g) + 2\,\ce{H2O}(l)$,\;
on donne (en \si{\kilo\joule\per\mol}) : $\dHf(\ce{CH4})=-75$, $\dHf(\ce{CO2})=-394$, $\dHf(\ce{H2O},l)=-286$.
\begin{enumerate}[label=\alph*),leftmargin=2em,topsep=2pt,itemsep=2pt]
  \item Écrire la loi de Hess pour cette réaction.
  \item Calculer $\dHr$. La réaction est-elle exo ou endothermique ?
\end{enumerate}
\end{exercice}

\begin{exercice}[inverser une réaction]
La synthèse \;$2\,\ce{SO2}(g) + \ce{O2}(g) \rightarrow 2\,\ce{SO3}(g)$\; a pour enthalpie
$\dH = \SI{-200}{\kilo\joule}$.
\begin{enumerate}[label=\alph*),leftmargin=2em,topsep=2pt,itemsep=2pt]
  \item Que vaut $\dH$ de la réaction inverse \;$2\,\ce{SO3}(g) \rightarrow 2\,\ce{SO2}(g) + \ce{O2}(g)$ ?
  \item La réaction inverse est-elle exo ou endothermique ?
\end{enumerate}
\end{exercice}

\begin{exercice}[effet des coefficients]
La formation de l'eau liquide \;$\ce{H2}(g) + \tfrac12\,\ce{O2}(g) \rightarrow \ce{H2O}(l)$\; a
$\dH = \SI{-286}{\kilo\joule}$.
\begin{enumerate}[label=\alph*),leftmargin=2em,topsep=2pt,itemsep=2pt]
  \item Que vaut $\dH$ pour \;$2\,\ce{H2}(g) + \ce{O2}(g) \rightarrow 2\,\ce{H2O}(l)$ ?
  \item Que vaut $\dH$ pour \;$\tfrac12\,\ce{H2}(g) + \tfrac14\,\ce{O2}(g) \rightarrow \tfrac12\,\ce{H2O}(l)$ ?
  \item Que représente la valeur $\SI{-286}{\kilo\joule\per\mol}$ vis-à-vis de $\ce{H2O}(l)$ ?
\end{enumerate}
\end{exercice}

\end{document}
